	\vspace{12pt}
	\section{Erläuterungen}
		Diese Prüfungsordnung ist, wie bereits erläutert, nach den Vorgaben des \href{https://www.karate.de/}{DKV} aufgebaut. Im Folgenden werden exemplarische, karatespezifische Vorgaben erläutert.
		\subsection{Techniken}
			\subsubsection{Angriffstechniken der oberen Extremitäten}
				\begin{tabularx}{\textwidth}{p{100pt}X}
					Zuki			&	Schlag/Stoß mit der geschlossenen Faust\\
					Teisho-Uchi		&	Angriffstechnik mit dem Handballen\\
					Uraken-Uchi		&	Angriff mit dem Faustrücken (schnappende Bewegung)\\
					Empi-Age-Uchi	& Ellbogenschlag von unten nach oben\\
				\end{tabularx}
			\subsubsection{Angriffstechniken der unteren Extremitäten}
				\begin{tabularx}{\textwidth}{p{100pt}X}
					Mae-Geri		&	Fußtritt nach vorne mit dem Fußballen\\
					Mawashi-Geri	&	Halbkreistritt - im Kihon-Ido ausschließlich in den mittleren Bereich und mit dem Fußballen, als Chisai ohne Drehung des Standbeins\\
				\end{tabularx}
			\subsubsection{Abwehrtechniken}
				\begin{tabularx}{\textwidth}{p{100pt}X}
					Age-Uke				&	Abwehr im oberen Bereich (Kopf)\\
					Yoko-Uke			&	Abwehr im mittleren Bereich (Bauch), der Arm arbeitet von innen nach außen\\
					Harai-Otoshi-Uke	& kreisrunde Abwehr im unteren Bereich (unter der Gürtellinie)\\
					Soto-Uke			&	Abwehr im mittleren Bereich (Bauch) - der Arm arbeitet von außen nach innen\\
					Kake-Uke			&	Hakenabwehr mit der Handinnenseite (offene Hand)\\
					Mawashi-Uke			&	kreisförmige Abwehr mit beiden Armen
				\end{tabularx}
		\subsection{Stände/Stellungen}
			\begin{tabularx}{\textwidth}{p{100pt}X}
				Zenkutsu-Dachi		&	lange Vorwärtsstellung - das vordere Knie ist angewinkelt, das hintere Bein weitestgehend gestreckt\\
				Sanchin-Dachi		&	kurze Vorwärtsstellung - der vordere Fuß ist leicht nach innen gerichtet, beide Knie zeigen eine angewinkelte Tendenz nach innen\\
				Shiko-Dachi (45°)	& Breiter Stand - beide Knie fest nach außen gedrückt, Oberkörper gerade\\
				Neko-Ashi-Dachi		&	\frqq Katzenfußstellung\flqq~ - der vordere Fuß steht auf dem Fußballen, Belastung zu ca. 70\% auf dem hinteren Bein, der Gesäßmuskel wird dabei nach unten gedrückt\\
				Suri-Ashi-Dachi		&	Parallele Fußhaltung - der vordere Fuß ca. 50\% vor dem anderen Fuß, Schulterbreit\\
			\end{tabularx}
	\section{Anmerkungen zu dieser Prüfungsordnung}
		Form, Layout und Design dieser Prüfungsordnung orientieren sich an der Farbgebung unseres Vereins, der \href{https://www.sglangenfeld.de/}{SGL}. Das Logo unserer Abteilung wurde von Steffi Gans erstellt und an die aktuelle Farbgebung angepasst. Als Basis für diese Neufassung der Prüfungsordnung diente die \frqq alte\flqq\,Prüfungsordnung, welche ebenfalls durch Steffi Gans erstellt wurde. Hier erfolgten Anpassungen im Bereich der karatespezifischen Anforderungen auf Basis der aktuellen Fassung der Prüfungsordnung des \href{https://www.karate-gkd.de/}{GKD}, sowie die Erweiterung um die Anforderungen zur Prüfung zum 1.\,Dan. Die einzelnen \frqq Prüfungskästen\flqq \,sind ohne den grünen Seitenrahmen gesetzt worden, um diese als CheatSheet besser ausdrucken zu können. Diese Prüfungsordnung erhebt keinen Anspruch auf Allgemeingültigkeit oder Vollständigkeit, sie dient zunächst der Orientierung.
	\section{Verwendete Literatur und Internetquellen}
		\printbibliography[heading=none]